\chapter{Idempotent Normalization and Realization}\label{chap:6}

\section*{Overview of the Chapter}

The question of this chapter is \textbf{what a comparison that includes a normalization
retains, and when one can return to the original}.
Chapter~\ref{chap:5} built two routes from the transport of structure and the selection of
the input, and obtained the invertible canonical comparison $\alpha$ between them.
It further incorporated a normalization $E$ chosen according to the diagnosis, and
constructed the generated comparison $\beta=E\alpha$.
In this chapter we examine how this last step changes the meaning of the comparison.

For example, consider a situation in which a design split into an order service and an
inventory service is brought together into a common design model.
We want to keep the operations of cancellation request, release of an inventory
reservation, and cancellation completion together with their dependencies, and to align the
class structure inside the services and the auxiliary data representations with a chosen
standard representation.
If we can check that the operations, Laws, and evaluation values that are kept are
unchanged by this processing, then the two design proposals can be compared on the standard
representation.
Whether the original class structure can be recovered from that comparison, however, must
be checked separately.

A normalization is idempotent when applying the alignment to the standard representation a
second time does not change the result.
Idempotence is the property that ``standardization is stable under repetition,'' and it
differs from invertibility, the property that ``the original representation can be
recovered.''
In this chapter we explain this difference at three stages: the map of objects, the
structure-preserving morphism, and the image after normalization.

\begin{xltabular}{\linewidth}{@{}LLL@{}}
\toprule
Question & Construction & Outcome \\
\midrule
\endhead
Which objects remain after normalization? & The object selected for each configuration & Correspondence between fixed points and configurations, unique factorization of readouts \\
\addlinespace[3pt]
Can normalization be repeated including operations and Laws? & Endomorphisms of cores and full geometries & Idempotence of the entire morphism \\
\addlinespace[3pt]
When is a comparison with a normalization added invertible? & Factorization into the canonical comparison and an idempotent factor & Equivalence with the normalization factor being the identity \\
\addlinespace[3pt]
How is a noninvertible comparison to be read between the pieces of information that remain? & The image in the idempotent completion & Isomorphism between the images given by the same generated comparison \\
\addlinespace[3pt]
Can the order of a normalization and an arbitrary design change be interchanged? & The normalization functor and the condition of naturality & A one-sided equation that always holds, and a counterexample from additional information on operations \\
\addlinespace[3pt]
Can a comparison itself be made an object of change? & The arrow category and idempotent completion & A category for classifying comparison-preserving changes in Chapter~\ref{chap:7} \\
\bottomrule
\end{xltabular}

In \S\S\ref{sec:6.1}--\ref{sec:6.4} we analyze a single generated comparison and proceed as
far as the classification of invertibility for full geometries.
In \S\S\ref{sec:6.5}--\ref{sec:6.7} we construct the changes between objects and the
category of comparisons, and in \S\ref{sec:6.8} we return to the semantics of lenses and
protocols.

\section{Selecting one representation for each configuration}\label{sec:6.1}

The architecture object of Chapter~\ref{chap:1} was an object obtained by adding structure
data and a quantity to a configuration.
What remains when the objects with the same configuration are aligned to a single one?
We first answer this question at the level of the map of objects, before operations and
Laws are added.

\begin{construction}[Projection to configurations and selection]\label{cons:6.1}

We fix the vocabulary of Atoms and write $\mathcal{A}$ for the totality of architecture
objects and $\mathcal{C}$ for the totality of configurations. For a core $P$ we set
\begin{equation*}
\pi:\mathcal{A}\longrightarrow\mathcal{C},\quad \pi(A)=C_A,
\qquad
s_P:\mathcal{C}\longrightarrow\mathcal{A},\quad s_P(C)=\mathrm{Form}_P(C)
\tag{6.1}\label{eq:6.1}
\end{equation*}
Here $\mathcal{A},\mathcal{C}$ are collections of objects, and \eqref{eq:6.1} defines maps
between them.
By the conditions on object formation, $\pi s_P=1_{\mathcal{C}}$.
Hence the normalization of Definition~\ref{def:5.39} satisfies
\begin{equation*}
n_P=s_P\pi,\qquad \pi n_P=\pi,\qquad n_P^2=n_P
\tag{6.2}\label{eq:6.2}
\end{equation*}
The last equation is the computation $s_P\pi s_P\pi=s_P\pi$.

\end{construction}

\begin{proposition}[Fixed points and configurations]\label{prop:6.2}

Between the collection of fixed points $\mathrm{Fix}(n_P)=\{A\mid n_P(A)=A\}$ and
$\mathcal{C}$ there are mutually inverse correspondences
\begin{equation*}
\begin{aligned}
\pi:\mathrm{Fix}(n_P)&\longrightarrow\mathcal{C},& A&\longmapsto C_A,\\
s_P:\mathcal{C}&\longrightarrow\mathrm{Fix}(n_P),& C&\longmapsto s_P(C)
\end{aligned}
\tag{6.3}\label{eq:6.3}
\end{equation*}
Moreover, if we define $A\sim B$ by $C_A=C_B$, then
$\mathcal{A}/{\sim}\cong\mathcal{C}$.

\end{proposition}

\begin{proof}

Since $n_Ps_P=s_P$, the object $s_P(C)$ is a fixed point.
The equation $\pi s_P=1$ and the equation $s_P\pi(A)=A$ at fixed points give a two-sided
inverse.
For the quotient, $[A]\mapsto C_A$ is well defined, and the inverse is $C\mapsto[s_P(C)]$.
The configurations of an object and of its selected object are equal, so this is also a
two-sided inverse.

\end{proof}

\begin{theorem}[Readouts unchanged by normalization]\label{thm:6.3}

For any set of values $Y$ and any map $f:\mathcal{A}\to Y$,
\begin{equation*}
fn_P=f
\quad\Longleftrightarrow\quad
\exists!\,\widetilde f:\mathcal{C}\to Y,
\qquad f=\widetilde f\pi.
\tag{6.4}\label{eq:6.4}
\end{equation*}

\end{theorem}

\begin{proof}

When the left-hand side holds, setting $\widetilde f=fs_P$ gives
$\widetilde f\pi=fn_P=f$.
If another factor $g$ satisfies $g\pi=f$, composing with $s_P$ on the right gives
$g=fs_P$. Conversely, if $f=\widetilde f\pi$, then
$fn_P=f$ by $\pi n_P=\pi$.

\end{proof}

This theorem characterizes the information that a normalization retains, from the point of
view of readouts.
A value that is unchanged before and after normalization can be read uniquely from the
configuration alone.
This theorem can also be applied to the invariance of residuals, invariants, and signatures
required in Definition~\ref{def:5.39}.
In that case we read with the arguments other than the object, such as the context, the
equation index, and the axis, held fixed.

\begin{example}[A finite model with two representations]\label{ex:6.4}

We take the set of configurations to be $\{c_0,c_1\}$, the tags of the representations
to be $\{0,1\}$, the objects to be $(c_i,t)$, and the selected objects to be $(c_i,0)$.
\begin{equation*}
n(c_i,t)=(c_i,0).
\tag{6.5}\label{eq:6.5}
\end{equation*}
Two of the four objects are fixed points.
The readout $f(c_i,t)=i$ is unchanged by the normalization and can be read from the
configuration.
The readout $g(c_i,t)=t$, on the other hand, distinguishes the original representations, so
the same value cannot be recovered after the normalization.
In the software example, one may take the selected operations and dependencies to
correspond to the former, and which auxiliary data represent that structure to correspond
to the latter.

\end{example}

\section{Lifting object normalization to a structure-preserving idempotent}\label{sec:6.2}

Even if a map of objects is idempotent, that alone does not make it a morphism preserving
operations and Laws.
The condition $\mathrm{Ad}(P)$ of Definition~\ref{def:5.39} supplies the preservation laws
needed for this.
In this section we show that the entire morphism built from that condition is idempotent.

\begin{theorem}[Idempotence of the normalization of a core]\label{thm:6.5}

For a core $P$ satisfying $\mathrm{Ad}(P)$, the endomorphism
$N_P:P\to P$ of Definition~\ref{def:5.39} satisfies
\begin{equation*}
N_P^2=N_P,\qquad q(N_P)=1_{q(P)}
\tag{6.6}\label{eq:6.6}
\end{equation*}
Here $q:E_{\mathrm{core}}\to B$ is the projection to the base.

\end{theorem}

\begin{proof}

For the object map we use \eqref{eq:6.2}.
The maps of Atoms, the base, contexts, equation indices, invariant indices, and signature
axes are identities, so composing them twice gives the same maps.
The correspondence of the rings of the equations and the correspondence of the value ranges
of the signatures are identities as well, and preservation of the residuals is by
$\mathrm{Ad}(P)$.

For the operations we use the identification along the equality of the types
$\mathrm{Op}_P(A,B)$ and $\mathrm{Op}_P(n_P(A),n_P(B))$.
Composing the identification twice is equal to the identification taken once, after the
endpoints have been aligned by $n_P^2=n_P$.
Its configuration map is the same as well, by the operation-preservation condition of
$\mathrm{Ad}(P)$.
This shows that all the data of a morphism in Definition~\ref{def:1.28} agree.
The base morphism is the identity by construction.

\end{proof}

\begin{lemma}[Naturality for the object map]\label{lem:6.6}

Let $H$ be the object map of an arbitrary morphism of cores $h:P\to Q$. Then
\begin{equation*}
Hn_P=n_QH.
\tag{6.7}\label{eq:6.7}
\end{equation*}
Neither $\mathrm{Ad}(P)$ nor $\mathrm{Ad}(Q)$ is needed for this conclusion.

\end{lemma}

\begin{proof}

Let $e$ be the bijection of Atoms of $h$. By the preservation of object formation and of
configurations in Definition~\ref{def:1.28}, we obtain
\[
H(n_P(A))
=H(\mathrm{Form}_P(C_A))
=\mathrm{Form}_Q(e_*C_A)
=\mathrm{Form}_Q(C_{H(A)})
=n_Q(H(A))
\]

\end{proof}

Equation~\eqref{eq:6.7} is an equality about the correspondence of objects.
The condition $hN_P=N_Qh$, which includes the correspondence of the operations as well, is examined in \S\ref{sec:6.6}.

\begin{proposition}[Distinct objects over the same configuration]\label{prop:6.7}

Over any configuration $C$ there exist at least two distinct architecture objects.
Hence, for every core $P$, the map $n_P$ is not injective, and if $\mathrm{Ad}(P)$ holds,
then $N_P$ is not an isomorphism.

\end{proposition}

\begin{proof}

Take the types of the structure data to be the one-element set $\{*\}$ and the two-element
set $\{0,1\}$, respectively, and the values of the structure data to be $*$ and $0$,
respectively.
For the quantities, in both objects we take the type to be a one-element set and the value
to be its unique element.
These are objects with the same $C$, but they are not equal, since the cardinalities of the
types of their structure data differ.
Both are normalized to $\mathrm{Form}_P(C)$, so $n_P$ is not injective.
If a morphism of cores had an inverse, its object map would have a two-sided inverse as
well, contradicting the noninjectivity of $n_P$. Hence $N_P$ is not an isomorphism.

\end{proof}

What we treat here is not only the selected representations but the totality of
architecture objects of Definition~\ref{def:1.6}.
This collection contains objects that represent the same configuration with different
structure data.
Collecting that distinction into a single selected object is what produces the
noninvertibility of $N_P$.

\begin{theorem}[Normalization does not split within the category of cores]\label{thm:6.8}

Suppose $\mathrm{Ad}(P)$ holds.
There exist no core $Q$ and morphisms of cores $r:P\to Q$ and $i:Q\to P$ satisfying
\begin{equation*}
ri=1_Q,\qquad ir=N_P
\tag{6.8}\label{eq:6.8}
\end{equation*}

\end{theorem}

\begin{proof}

Suppose they exist, and let the object maps of $r,i$ be $R,I$, respectively.
Since $RI=1$, the map $I$ is injective.
By Lemma~\ref{lem:6.6} and $IR=n_P$, for every object $A$ we obtain
\[
I(n_Q(A))=n_P(I(A))=IR(I(A))=I(A)
\]
Injectivity gives $n_Q(A)=A$, contradicting Proposition~\ref{prop:6.7}.

\end{proof}

For the collections of objects, \eqref{eq:6.2} gave a splitting through the configurations.
A core, however, is not the set of configurations itself: it carries the readings of object
formation, operations, and Laws for all architecture objects. Merely extracting the set of
configurations does not yield the intermediate core of \eqref{eq:6.8}.
To treat the image of a normalization as an object of a category, in the next section we
extend the objects and the morphisms.

\section{Making the remaining information into objects by idempotent completion}\label{sec:6.3}

For an idempotent map on a set, the set of its fixed points can be extracted as its image.
The idempotent completion is what makes it possible to treat objects playing the same role
in a general category.
By this construction, the information that the original comparison loses and the
correspondence preserved between the images can be expressed by the same equations.

\begin{definition}[Karoubi envelope]\label{def:6.9}

An object of the idempotent completion $\mathrm{Kar}(E)$ of a category $E$ is a pair
$(X,e)$ of an object $X\in E$ and an idempotent $e:X\to X$.
A morphism $a:(X,e)\to(Y,d)$ is a morphism $a:X\to Y$ of $E$ satisfying
\begin{equation*}
dae=a
\quad\Longleftrightarrow\quad
da=a=ae
\tag{6.9}\label{eq:6.9}
\end{equation*}
For the composition we use the composition of $E$, and we take the identity morphism of
$(X,e)$ to be $e$.
These form a category, by the category laws of $E$ and \eqref{eq:6.9}.

\end{definition}

\begin{construction}[Embedding of the original category and splitting of idempotents]\label{cons:6.10}

Define $J:E\to\mathrm{Kar}(E)$ by $J(X)=(X,1_X)$ and $J(f)=f$.
It is fully faithful. For any $(X,e)$ there are two morphisms, each of which is $e$ as a
morphism of the original category,
\begin{equation*}
J(X)\xrightarrow{\ r_X\ }(X,e)
\xrightarrow{\ i_X\ }J(X)
\tag{6.10}\label{eq:6.10}
\end{equation*}
and
\begin{equation*}
r_Xi_X=1_{(X,e)},\qquad i_Xr_X=J(e).
\tag{6.11}\label{eq:6.11}
\end{equation*}

\end{construction}

\begin{proof}

Between objects $(X,1_X)$, condition \eqref{eq:6.9} holds for every morphism, so $J$ is
fully faithful.
Checking \eqref{eq:6.9} for $r_X,i_X$, we find that both amount to $e^2=e$.
The underlying morphism of both composites is $e$, which is the identity on $(X,e)$ and is
$J(e)$ on $J(X)$.

\end{proof}

\begin{proposition}[Realization of the image of an idempotent]\label{prop:6.11}

In $\mathrm{Kar}(E)$ every idempotent splits.
Moreover, when there is a splitting in $E$ satisfying $ri=1_Z$ and $ir=e$, we have
$(X,e)\cong J(Z)$.

\end{proposition}

\begin{proof}

An idempotent morphism $a:(X,e)\to(X,e)$ satisfies $a^2=a$ and $ea=a=ae$.
Forming the object $(X,a)$ and taking the underlying morphism in both directions to be $a$
gives its splitting.
For the second claim, we regard the given $r,i$ as morphisms between $(X,e)$ and $J(Z)$.
The equations $re=r$ and $ei=i$ and the two-sided inverse equations follow from the
equations of the splitting.

\end{proof}

Hence $(P,N_P)$ becomes a new object that handles the information retained by the
normalization of a core.
Theorem~\ref{thm:6.8} shows that this is not an object obtained by splitting the same
normalization inside $E_{\mathrm{core}}$.
In the idempotent completion, taking the identity morphism to be $N_P$ specifies the
morphisms appropriate to that image.

\begin{theorem}[Normalization factorization of a comparison and isomorphism of images]\label{thm:6.12}

In an arbitrary category $E$, take an isomorphism $\alpha:X\xrightarrow{\sim}Y$ and an
idempotent $d:Y\to Y$, and set
\begin{equation*}
\beta=d\alpha,\qquad e=\alpha^{-1}d\alpha
\tag{6.12}\label{eq:6.12}
\end{equation*}
Then $e^2=e$, and the following hold.
\begin{equation*}
\begin{gathered}
\beta\text{ is invertible in }E\text{}
\quad\Longleftrightarrow\quad d=1_Y,\\
\beta:(X,e)\xrightarrow{\sim}(Y,d)
\quad\text{in }\mathrm{Kar}(E).
\end{gathered}
\tag{6.13}\label{eq:6.13}
\end{equation*}
The underlying morphism of the inverse between the images is $\gamma=\alpha^{-1}d$.

\end{theorem}

\begin{proof}

The computation with conjugation gives $e^2=\alpha^{-1}d^2\alpha=e$.
Since $\alpha$ is invertible, invertibility of $\beta$ is equivalent to invertibility of
$d$.
For an invertible idempotent, cancelling one $d$ in $d^2=d$ gives $d=1_Y$.
The converse is clear.

Next we check $d\beta e=\beta$ and $e\gamma d=\gamma$ using $d^2=d$.
The two composites are
\begin{equation*}
\gamma\beta=e=1_{(X,e)},\qquad
\beta\gamma=d=1_{(Y,d)}
\tag{6.14}\label{eq:6.14}
\end{equation*}
so $\gamma$ is the inverse morphism in $\mathrm{Kar}(E)$.

\end{proof}

What becomes isomorphic in this result are the images that remain after the normalization.
The morphism used for the isomorphism of the images is the same $\beta$ that was not
invertible in the original category, and what has changed are the objects at the two ends
and their identity morphisms.

\begin{example}[From four points to an image with two points]\label{ex:6.13}

Take the sets $X=\{a,b,c,d\}$ and $Y=\{0,1,2,3\}$, and give the isomorphism $\alpha$ and
the normalization $d_Y$ by the following table.

\begin{xltabular}{\linewidth}{@{}LLLL@{}}
\toprule
$x$ & $\alpha(x)$ & $\beta(x)=d_Y(\alpha(x))$ & $e_X(x)$ \\
\midrule
\endhead
$a$ & $1$ & $0$ & $c$ \\
\addlinespace[3pt]
$b$ & $3$ & $2$ & $d$ \\
\addlinespace[3pt]
$c$ & $0$ & $0$ & $c$ \\
\addlinespace[3pt]
$d$ & $2$ & $2$ & $d$ \\
\bottomrule
\end{xltabular}

Here $d_Y(0)=d_Y(1)=0$, $d_Y(2)=d_Y(3)=2$, and
$e_X=\alpha^{-1}d_Y\alpha$.
Since $\beta$ identifies $a,c$ with each other and $b,d$ with each other, it is not
invertible between the four points.
Between the sets of fixed points $\{c,d\}$ and $\{0,2\}$, on the other hand, it is a
bijection.
Equation~\eqref{eq:6.14} expresses this correspondence of images in a general category.

\end{example}

\section{Classifying the generated comparisons of full geometries}\label{sec:6.4}

In the generated comparison of Chapter~\ref{chap:5}, whether to normalize was determined by
the diagnosis and the preservation conditions, and that endomorphism was carried along the
actual route.
To apply Theorem~\ref{thm:6.12} to this construction, we verify idempotence for morphisms
that include the covers, the overlaps, the coefficients, and the local data, in addition to
the operations and Laws of a core.

\begin{theorem}[Idempotent normalization of a full geometry]\label{thm:6.14}

For a full geometry $G$ whose core $P=p(G)$ satisfies $\mathrm{Ad}(P)$, the endomorphism
$N_G:G\to G$ of Definition~\ref{def:5.39} satisfies
\begin{equation*}
N_G^2=N_G,\qquad p(N_G)=N_P,\qquad
qp(N_G)=1_{qp(G)}
\tag{6.15}\label{eq:6.15}
\end{equation*}
The map to the coefficient ring is the identity.

\end{theorem}

\begin{proof}

Idempotence of the core component is by Theorem~\ref{thm:6.5}.
In the remaining components of a morphism of geometries, the maps of contexts, Atoms, and
axes are identities, the coverage requirements are preserved as they are, and the
comparison of overlaps is the identity isomorphism.
The coefficient homomorphism and the comparison maps of supports, axes, and observables are
identities as well, so the twofold composite equals the single one.
The raw system is preserved by the identity reindexing and the identity change of
coefficients.
Hence all the components of Definition~\ref{def:1.30} agree. The equations for the
projections follow from the construction.

\end{proof}

\begin{lemma}[Interchange of generated transport and normalization]\label{lem:6.15}

The canonical transport of cores along an arbitrary exact base morphism $\sigma:b\to b'$,
and the pullback of cores by the reindexing in the opposite direction, preserve
$\mathrm{Ad}$.
The canonical forward transport of geometries also preserves the same condition on cores.

Assume further that Atom equality is decidable, and let $\sigma$ be the decoding of a
morphism of finite codes of Definition~\ref{def:5.14}.
The canonical pullback of full geometries for this input also preserves $\mathrm{Ad}$, and
for the two functors of transport and pullback
\begin{equation*}
\sigma_!(N_G)=N_{\sigma_!G},\qquad
\sigma^*(N_H)=N_{\sigma^*H}
\tag{6.16}\label{eq:6.16}
\end{equation*}
hold. Here the cores of $G,H$ satisfy $\mathrm{Ad}$.

For the actual lifts of cores along an arbitrary exact morphism,
\begin{equation*}
N_{\sigma_!P}\iota_{\sigma,P}=\iota_{\sigma,P}N_P,
\qquad
N_Q\kappa_{\sigma,Q}=\kappa_{\sigma,Q}N_{\sigma^*Q}
\tag{6.17}\label{eq:6.17}
\end{equation*}
hold. The same interchange equations hold for the forward lifts of geometries.
For the pullback lifts of geometries we use the morphism realized by the finite code above.

\end{lemma}

\begin{proof}

Write $e$ for the Atom map of the base morphism and $T_e$ for the reindexing of objects.
From the definition of the transported object formation we obtain
\begin{equation*}
n_{\sigma_!P}(T_eA)=T_e(n_P(A))
\tag{6.18}\label{eq:6.18}
\end{equation*}
Since $T_e$ has $T_{e^{-1}}$ as its inverse,
the preservation laws on the target of the transport can be brought back to the
preservation laws before the transport.
Carrying the invariance of residuals, invariants, and signatures, the equality of the types
of operations, and the commutativity of the configuration maps along this correspondence,
we see that $\mathrm{Ad}$ is preserved. For the pullback of cores we use $e^{-1}$.

The operation map of a generated lift is built from the reindexing and the identification
by the equality of the types at the endpoints.
Moving the normalization before or after it, aligning the endpoints by \eqref{eq:6.18}
yields the same operation.
The other core components agree as well, and \eqref{eq:6.17} holds.
For the forward lifts of geometries and the pullback lifts from a finite code, composing
the identity local comparisons of Theorem~\ref{thm:6.14} gives the interchange equations
for all the components.
Finally, using the universal property of each lift for this finite-code input, the image of
the normalization morphism in \eqref{eq:6.16} is uniquely determined.

\end{proof}

Naturality of the object map for a general morphism of cores is by Lemma~\ref{lem:6.6}.
The interchange condition for morphisms that have their own operation map beyond the
objects is examined in \S\ref{sec:6.6}.

\begin{theorem}[Normalization factorization of the generated comparison and classification of invertibility]\label{thm:6.16}

We fix the common input of Construction~\ref{cons:5.37}, namely the square decoded from
finite codes, the core at each vertex and the strong edges, the conditions on the end of
each face and on the two paths in the base, and the specified comparisons.
From a face $z$ of the same diagram, a cochain $\omega$, the core $P_z$ at its end, a
commutative coefficient ring $k$, and a full geometry $G_z$, we form the comparison of
Construction~\ref{cons:5.40}.
We abbreviate the ends of the two routes and the canonical comparison as
\begin{equation*}
X=\overline D(G_z),\qquad Y=\overline V(G_z),\qquad
\overline\alpha=\overline\alpha_z:X\xrightarrow{\sim}Y
\tag{6.19}\label{eq:6.19}
\end{equation*}
The condition for selecting the normalization and the generated morphisms are taken to be
\begin{equation*}
\begin{aligned}
\chi&\Longleftrightarrow
   (\omega_z\ne1)\ \text{and}\ \mathrm{Ad}(P_z),\\
\overline d&=\overline V(\overline S_{z,\omega}),\qquad
\overline\beta=\overline d\,\overline\alpha,\qquad
\overline e=\overline\alpha^{-1}\overline d\,\overline\alpha
\end{aligned}
\tag{6.20}\label{eq:6.20}
\end{equation*}
Then $\overline d^2=\overline d$ and $\overline e^2=\overline e$, and
\begin{equation*}
\overline\beta\text{ is invertible in }E_{\mathrm{geom}}\text{}
\quad\Longleftrightarrow\quad
\overline d=1_Y
\quad\Longleftrightarrow\quad
\neg\chi.
\tag{6.21}\label{eq:6.21}
\end{equation*}
Moreover, the same morphism $\overline\beta$ gives an isomorphism
\begin{equation*}
(X,\overline e)\xrightarrow{\ \overline\beta\ }(Y,\overline d)
\quad\text{in }\mathrm{Kar}(E_{\mathrm{geom}})
\tag{6.22}\label{eq:6.22}
\end{equation*}
and the underlying morphism of its inverse is $\overline\alpha^{-1}\overline d$.
When $\chi$ holds, the cores at both ends satisfy $\mathrm{Ad}$ as well, and
\begin{equation*}
\overline d=N_Y,\qquad \overline e=N_X.
\tag{6.23}\label{eq:6.23}
\end{equation*}
Hence, when the normalization is selected, the same generated comparison matches the
normalization images $(X,N_X)$ and $(Y,N_Y)$, determined at the respective ends of the
two routes, by an isomorphism.
The idempotent at the start was defined by conjugation with the comparison, and we see
further that it agrees with the normalization $N_X$ of the start itself.
When $\chi$ does not hold, both idempotents are identities.

\end{theorem}

\begin{proof}

The endomorphism on the source in Construction~\ref{cons:5.40} is $N_{G_z}$ if $\chi$ holds
and the identity otherwise.
By Theorem~\ref{thm:6.14} and functoriality, its image $\overline d$ is idempotent.
Since $\overline\alpha$ of Construction~\ref{cons:5.37} is an isomorphism,
Theorem~\ref{thm:6.12} gives the first equivalence of \eqref{eq:6.21}, the idempotence of
$\overline e$, and \eqref{eq:6.22}.

Suppose $\chi$ holds. Applying Lemma~\ref{lem:6.15} to the transports and pullbacks along
the four edges gives $\mathrm{Ad}$ at both ends and $\overline d=N_Y$.
Furthermore, into the geometric version of \eqref{eq:5.18}, which characterizes the
canonical comparison,
\begin{equation*}
\kappa_{\sigma_2,(\sigma_1)_!G_z}\,
\overline\alpha\,\iota_{\pi_2,\pi_1^*G_z}
=\iota_{\sigma_1,G_z}\,\kappa_{\pi_1,G_z}
\tag{6.24}\label{eq:6.24}
\end{equation*}
we substitute \eqref{eq:6.17} for each lift.
After composing with the lifts on both outer sides, $N_Y\overline\alpha$ and
$\overline\alpha N_X$ become the same morphism.
Using in turn the uniqueness of each lift for cores and for geometries gives
$N_Y\overline\alpha=\overline\alpha N_X$, which is the second equation of \eqref{eq:6.23}.

By Proposition~\ref{prop:6.7}, $N_{p(Y)}$ is not the identity.
Since $p(N_Y)=N_{p(Y)}$, we have $\overline d=N_Y\ne1_Y$.
If, on the other hand, $\neg\chi$, then both the endomorphism on the source and its image
are identities, and the conjugated $\overline e$ is the identity as well. This gives the
remaining equivalence of \eqref{eq:6.21}.

\end{proof}

Equation~\eqref{eq:6.21} is a classification for the generation rule chosen in
Construction~\ref{cons:5.40}.
Even when the diagnosis is not the identity, this rule selects the identity morphism unless
the preservation conditions for the normalization are satisfied.
Moreover, selecting a normalization and removing the value of the diagnosis are
separate statements.
What this theorem clarifies is the invertibility and the image of the comparison generated
from that selection.

\begin{proposition}[Projection to the comparison of cores]\label{prop:6.17}

For the input of Theorem~\ref{thm:6.16} we write $\alpha=\alpha_{P_z}$,
$\beta=\beta_{z,\omega}$, and $e=\alpha^{-1}E_{z,\omega}\alpha$.
Under the endpoint isomorphisms $j_D:pX\xrightarrow{\sim}DP_z$ and
$j_V:pY\xrightarrow{\sim}VP_z$ of Proposition~\ref{prop:5.38},
\begin{equation*}
\begin{aligned}
j_Vp(\overline\alpha)&=\alpha j_D,&
j_Vp(\overline d)&=E_{z,\omega}j_V,\\
j_Dp(\overline e)&=e j_D,&
j_Vp(\overline\beta)&=\beta j_D.
\end{aligned}
\tag{6.25}\label{eq:6.25}
\end{equation*}
Hence, projecting the isomorphism of images of geometries to the core gives the isomorphism
of images constructed at the core.

\end{proposition}

\begin{proof}

The first two equations and the last one are \eqref{eq:5.42} and \eqref{eq:5.46}.
The remaining equation is obtained by projecting
$\overline e=\overline\alpha^{-1}\overline d\,\overline\alpha$, using that $p$ preserves
composites and inverses.
In the idempotent completion, the underlying morphisms of the isomorphisms that align the
endpoints are $e j_D$ and $E_{z,\omega} j_V$, respectively.
By \eqref{eq:6.25} and idempotence, these define isomorphisms between the images and make
the square of the comparisons commute.

\end{proof}

\begin{example}[A normalization that collects the objects over the same configuration]\label{ex:6.18}

To see the action of a normalization, we define the data of a core as follows.

\begin{itemize}
\item \textbf{Atoms and sources.} We take one Atom and one source, let the extracted family
be that one-element set of Atoms, and take the source normalization to be the identity.
\item \textbf{Composition and object formation.} We set
$\mathrm{Comp}(F)=(F,\varnothing,\varnothing)$.
For every configuration we take the types of the structure data and of the quantity to be
one-element sets and their values to be the unique elements.
\item \textbf{Operations.} We set $\mathrm{Op}_P(A,B)=\mathrm{Hom}(C_A,C_B)$.
\item \textbf{Signatures.} We take two axes and let both values be the integer 0. Both axes
are taken to be selected axes.
\item \textbf{Contexts and local readouts.} A context is a pair $(U,T)$ of a subset $U$ of
the extracted family and a subset $T$ of the signature axes, and the morphisms are the
componentwise inclusions. The support reads each Atom of $U$ as it is, and all of $T$ can
be read as axes.
The observables form a one-element set that can be read, and for the restrictions we use
the inclusions on supports and axes and the identity on observables.
\item \textbf{Equations and invariants.} We take both index sets to be empty, the ring of
the equations to be the constant $\mathbb{Z}$, and the restrictions to be identities.
Since the indices are empty, the conditions on the residuals and the specification of the
detectors are empty.
\end{itemize}

This core $P$ satisfies $\mathrm{Ad}(P)$.
The normalization $N_P$ of Theorem~\ref{thm:6.5} sends the objects over the same
configuration to the selected object.
By Proposition~\ref{prop:6.7}, $N_P$ is not invertible.

For the full geometry we add the following data.

\begin{itemize}
\item \textbf{Overlaps and covers.} On the contexts above we take the overlaps to be the
componentwise intersections and the coverage requirements to be empty.
\item \textbf{Coefficients and the raw system.} We take the coefficients to be $\mathbb{Z}$,
and use a single raw coordinate $x$ for each context together with the relation $x^2-x$.
We take the kind of the coordinates to be semantic, the type of the local data to be a
one-element set, and all the restrictions to be identities.
\end{itemize}

The restrictions preserve the relations and respect identities and composition, so a full
geometry $G$ is obtained.
Since $N_G$ of Theorem~\ref{thm:6.14} projects to $N_P$, it is not invertible for the full
geometry either.
By Definition~\ref{def:6.9}, the identity morphism of $(G,N_G)$ is this $N_G$.

\end{example}

\section{Collecting normalized objects and changes into a single category}\label{sec:6.5}

Having constructed the image for a single comparison, we next treat several cores and the
changes between them at the same time.
The aim is to build the morphisms between the images consistently from the original
morphisms, rather than to choose them anew at each design change.

\begin{lemma}[One-sided absorption law]\label{lem:6.19}

Let $\mathcal{C}_{\mathrm{ad}}$ be the full subcategory of cores satisfying
$\mathrm{Ad}$. For every morphism $f:P\to Q$ of this subcategory,
\begin{equation*}
N_QfN_P=fN_P.
\tag{6.26}\label{eq:6.26}
\end{equation*}

\end{lemma}

\begin{proof}

Let $F$ be the object map. By Lemma~\ref{lem:6.6},
$n_QFn_P=Fn_P^2=Fn_P$.
An operation $o\in\mathrm{Op}_P(A,B)$ is first identified with one at the normalized
endpoints and is then sent by the operation map of $f$.
Its target objects $F(n_P(A)),F(n_P(B))$ are already fixed points by the equality of object
maps just obtained.
Hence the operation map of the final $N_Q$ is the identity identification on the same type.
For the components of the base, Atoms, contexts, equations, invariant indices, axes, and
value ranges, only an identity is composed at the end.
All the components of the morphisms are equal, so we obtain \eqref{eq:6.26}.

\end{proof}

\begin{construction}[The category of normalized cores]\label{cons:6.20}

For the objects of $\mathcal{C}_{\mathrm{nor}}$ we use the same cores as in
$\mathcal{C}_{\mathrm{ad}}$. Its morphisms and identity morphisms are defined by
\begin{equation*}
\begin{aligned}
\mathrm{Hom}_{\mathcal{C}_{\mathrm{nor}}}(P,Q)
 &=\{a:P\to Q\text{ in }\mathcal{C}_{\mathrm{ad}}\mid N_QaN_P=a\},\\
1_P^{\mathrm{nor}}&=N_P
\end{aligned}
\tag{6.27}\label{eq:6.27}
\end{equation*}
The composition is the composition of the original morphisms of cores.
By \eqref{eq:6.9}, this is the category whose morphisms are all the Karoubi morphisms between the objects
$(P,N_P)$.
Hence
\begin{equation*}
K:\mathcal{C}_{\mathrm{nor}}\longrightarrow\mathrm{Kar}(\mathcal{C}_{\mathrm{ad}}),
\qquad K(P)=(P,N_P),\quad K(a)=a
\tag{6.28}\label{eq:6.28}
\end{equation*}
is a fully faithful functor. Even though the objects are denoted by the same symbol $P$,
the identity morphisms and the sets of morphisms have changed.

\end{construction}

\begin{theorem}[Fullness of the normalization functor]\label{thm:6.21}

The following assignment defines a full functor.
\begin{equation*}
\begin{gathered}
\mathsf{N}_{\mathrm{core}}:\mathcal{C}_{\mathrm{ad}}\longrightarrow\mathcal{C}_{\mathrm{nor}},\\
\mathsf{N}_{\mathrm{core}}(P)=P,\qquad
\mathsf{N}_{\mathrm{core}}(f)=fN_P.
\end{gathered}
\tag{6.29}\label{eq:6.29}
\end{equation*}

\end{theorem}

\begin{proof}

By \eqref{eq:6.26} and $N_P^2=N_P$, the morphism $fN_P$ is a morphism of \eqref{eq:6.27}.
The identity morphism is sent to $1_PN_P=N_P=1_P^{\mathrm{nor}}$.
For $f:P\to Q$ and $g:Q\to R$,
\begin{equation*}
(gN_Q)(fN_P)=g(N_QfN_P)=gfN_P
\tag{6.30}\label{eq:6.30}
\end{equation*}
so the composition is preserved.
Every morphism $a:P\to Q$ in $\mathcal{C}_{\mathrm{nor}}$ satisfies $aN_P=a$ by
\eqref{eq:6.9}.
Choosing the underlying morphism of cores $a$ itself as a preimage, we have
$\mathsf{N}_{\mathrm{core}}(a)=a$.

\end{proof}

Fullness is the assertion that every morphism permitted after the normalization has an
original morphism of cores.
Distinct morphisms of cores can have the same image, and \S\ref{sec:6.6} gives an example.

\begin{proposition}[Extension to full geometries and compatibility with the projection]\label{prop:6.22}

Let $\mathcal{G}_{\mathrm{ad}}$ be the full subcategory of full geometries whose cores
satisfy $\mathrm{Ad}$. For every morphism $f:G\to H$,
\begin{equation*}
N_HfN_G=fN_G.
\tag{6.31}\label{eq:6.31}
\end{equation*}
Hence we obtain the category $\mathcal{G}_{\mathrm{nor}}$ whose morphisms are defined by
$N_HaN_G=a$, and the full functor $\mathsf{N}_{\mathrm{geom}}(f)=fN_G$.
Writing $p_{\mathrm{ad}},p_{\mathrm{nor}}$ for the functors induced by the projection to
cores,
\begin{equation*}
p_{\mathrm{nor}}\mathsf{N}_{\mathrm{geom}}
=\mathsf{N}_{\mathrm{core}}p_{\mathrm{ad}}.
\tag{6.32}\label{eq:6.32}
\end{equation*}
Moreover, the base and the coefficient homomorphism of a morphism after the normalization
are equal to those of the original morphism.

\end{proposition}

\begin{proof}

For the core component we use Lemma~\ref{lem:6.19}.
In the components of the geometry, the final $N_H$ composes identity maps onto the
coefficients, supports, axes, and observables, so the result equals the original $fN_G$.
This gives \eqref{eq:6.31}.
Functoriality and fullness are the same computation as in Theorem~\ref{thm:6.21}.
By $p(N_G)=N_{p(G)}$, equation \eqref{eq:6.32} holds both on objects and on morphisms.
For the base and the coefficients, the components of $N_G$ are identities, so composing
with them does not change the original morphism.

\end{proof}

\section{Conditions for interchanging the order of a normalization and a change}\label{sec:6.6}

That the normalization is a functor does not mean that the order of an arbitrary change and
the normalization can be interchanged.
What the former required was \eqref{eq:6.26}, and the latter is a stronger equation.
We examine in which component that difference appears.

\begin{proposition}[Naturality of the inclusion and the condition for the projection]\label{prop:6.23}

We use $K$ of Construction~\ref{cons:6.20} and the embedding $J$ obtained by applying
Construction~\ref{cons:6.10} to $\mathcal{C}_{\mathrm{ad}}$.
The inclusion $i_P:(P,N_P)\to J(P)$ whose underlying morphism is $N_P$ defines a natural
transformation
\begin{equation*}
i:K\mathsf{N}_{\mathrm{core}}\Longrightarrow J
\tag{6.33}\label{eq:6.33}
\end{equation*}
The morphism $r_P:J(P)\to(P,N_P)$ in the opposite direction satisfies
$r_Pi_P=1_{(P,N_P)}$ objectwise, and naturality of that family with respect to a morphism
$f:P\to Q$ is equivalent to
\begin{equation*}
fN_P=N_Qf
\tag{6.34}\label{eq:6.34}
\end{equation*}

\end{proposition}

\begin{proof}

The underlying morphisms of the two sides of the naturality of the inclusion are $fN_P$ and
$N_QfN_P$, which are equal by Lemma~\ref{lem:6.19}.
The objectwise splitting is by \eqref{eq:6.11}.
Naturality of the family in the opposite direction is the equation $(fN_P)N_P=N_Qf$, which
becomes \eqref{eq:6.34} on using $N_P^2=N_P$.

\end{proof}

\begin{proposition}[Deciding naturality by the correspondence of operations]\label{prop:6.24}

Suppose $\mathrm{Ad}(P)$ and $\mathrm{Ad}(Q)$ hold, and let $F$ be the object map of
$f:P\to Q$.
We write the operation map of the normalization as
\begin{equation*}
c^P_{A,B}:\mathrm{Op}_P(A,B)
\longrightarrow\mathrm{Op}_P(n_P(A),n_P(B))
\tag{6.35}\label{eq:6.35}
\end{equation*}
Then \eqref{eq:6.34} is equivalent to the statement that, for all $A,B,o$,
\begin{equation*}
f^{\mathrm{op}}_{n_P(A),n_P(B)}(c^P_{A,B}(o))
=c^Q_{F(A),F(B)}(f^{\mathrm{op}}_{A,B}(o))
\tag{6.36}\label{eq:6.36}
\end{equation*}
holds. The endpoints of the two sides are identified by Lemma~\ref{lem:6.6}.

\end{proposition}

\begin{proof}

Equation~\eqref{eq:6.36} is the equality of the operation components of two composite
morphisms.
The object components are equal by Lemma~\ref{lem:6.6}, and the components of the base,
Atoms, contexts, equations, invariant indices, axes, and value ranges are equal because the
identity maps of the normalization are composed either before or after.
Hence agreement of the operation components is equivalent to agreement of the entire
morphisms.

\end{proof}

In \eqref{eq:6.36} we compare not only the action of an operation on configurations but the
operation itself.
As in Example~\ref{ex:1.8} of Chapter~\ref{chap:1}, operations with the same action can
differ in their names or in additional information.

\begin{example}[The order changes according to the tag of an operation]\label{ex:6.25}

We replace the operations of the core of Example~\ref{ex:6.18} by
\begin{equation*}
\mathrm{Op}_{P^{\mathrm{tag}}}(A,B)
=\mathrm{Hom}(C_A,C_B)\times\{0,1\}
\tag{6.37}\label{eq:6.37}
\end{equation*}
For the action on configurations we use only the first component, and the other readings
are kept as they are.
The second component is a tag that distinguishes two operations with the same action.
This core also satisfies $\mathrm{Ad}$, and the normalization keeps the configuration map
and the tag as they are.

Over a configuration $C$, let $A_\bullet$ be an object whose type of structure data is a
two-element set.
The type of the structure data of its selected object $n(A_\bullet)$ is a one-element set,
so $n(A_\bullet)\ne A_\bullet$.
Keeping the objects, the base, and the other readings identical, we apply to the operations
alone
\begin{equation*}
f^{\mathrm{op}}_{A,B}(u,t)=
\begin{cases}
(u,1-t),&A=A_\bullet,\\
(u,t),&A\ne A_\bullet
\end{cases}
\tag{6.38}\label{eq:6.38}
\end{equation*}
Since the first component $u$ is unchanged, the commutative square for configurations in
Definition~\ref{def:1.28} is satisfied.
The other components are identities as well, and $f$ is an actual endomorphism of cores.
Applying it twice returns to the original, so $f^2=1$ as well.

Taking the start and the end to be $A_\bullet$ and the operation to be $(1_C,0)$, the
following difference arises.

\begin{xltabular}{\linewidth}{@{}LLL@{}}
\toprule
Order of application & Start and end after normalization & Final tag \\
\midrule
\endhead
Normalize and then apply $f$ & Both $n(A_\bullet)$ & 0 \\
\addlinespace[3pt]
Apply $f$ and then normalize & Both $n(A_\bullet)$ & 1 \\
\bottomrule
\end{xltabular}

Hence $fN\ne Nf$, and the family in the opposite direction in Proposition~\ref{prop:6.23}
is not natural.
The object map and the action on configurations agree along both routes.
What detects the difference of order is the tag kept in the operations.

Furthermore, all the targets of the normalization have one-element structure data, so none
of them is $A_\bullet$.
From this, $fN=N=1N$.
Two distinct morphisms $f\ne1$ are sent to the same morphism after the normalization, so
the full functor of Theorem~\ref{thm:6.21} is not faithful in general.

Read in terms of design changes, even when the correspondence of operations after the
alignment to the standard representation is the same, the specification of operations that
depends on the original representation can differ.
For example, for the same cancellation operation one may consider processing that changes
the label for auditing according to the original implementation variant.
If the distinction of implementation variants is collected first, the operations whose
labels are to be changed can no longer be distinguished.
When one wants to interchange the order before and after the normalization freely, a review
checks \eqref{eq:6.36} in addition to the correspondence of objects.

\end{example}

\section{The category of comparisons themselves and the three-level projections}\label{sec:6.7}

Chapter~\ref{chap:7} studies the changes that preserve a single comparison.
For this it is convenient to take the comparison morphisms themselves as objects and to
treat the changes at the two ends as morphisms.
The idempotent completion is compatible with this category of comparisons as well.

\begin{definition}[The arrow category and the category of comparisons]\label{def:6.26}

The arrow category $\mathrm{Arr}(E)$ of a category $E$ has the morphisms $c:X\to Y$ as its
objects.
A morphism from $c$ to $c':X'\to Y'$ is a pair of $a:X\to X'$ and $b:Y\to Y'$ satisfying
\begin{equation*}
bc=c'a
\tag{6.39}\label{eq:6.39}
\end{equation*}
Identities and composition are defined componentwise at the two ends.
We take the category that handles the comparisons after normalization to be
\begin{equation*}
\mathcal{M}(E)=\mathrm{Arr}(\mathrm{Kar}(E))
\tag{6.40}\label{eq:6.40}
\end{equation*}

\end{definition}

\begin{theorem}[Interchange of comparison and idempotent completion]\label{thm:6.27}

For every category $E$ there is a natural equivalence of categories
\begin{equation*}
\Phi_E:\mathrm{Kar}(\mathrm{Arr}(E))
\simeq\mathrm{Arr}(\mathrm{Kar}(E))
\tag{6.41}\label{eq:6.41}
\end{equation*}
Representing an object of the left-hand side by a morphism $c:X\to Y$ and an idempotent
commutative square $(e,d)$, that is, by $e^2=e$, $d^2=d$, and $dc=ce$, the functor $\Phi_E$
sends it to
\begin{equation*}
(X,e)\xrightarrow{\ dce\ }(Y,d)
\tag{6.42}\label{eq:6.42}
\end{equation*}
This correspondence is shown in Figure~\ref{fig:6.1}.

\end{theorem}

\begin{figure}[htbp]
\centering
% Figure 6.1 -- from an idempotent commutative square to a morphism between images.
% Same content as the Japanese manuscript's figure ../../figures/ch06-karoubi-arrow.svg, redrawn in TikZ.
\begin{tikzpicture}[
    x=1mm, y=1mm,
    arr/.style={-{Stealth[length=4.5pt]}, semithick},
    mapsto/.style={{Bar[width=5pt]}-{Stealth[length=4.5pt]}, semithick},
    panel/.style={font=\small},
    note/.style={font=\small}
  ]
  % Left: the idempotent commutative square.
  \node[panel] at (15,31) {Idempotent commutative square};
  \node (XT) at (0,22)  {$X$};
  \node (YT) at (30,22) {$Y$};
  \node (XB) at (0,0)   {$X$};
  \node (YB) at (30,0)  {$Y$};
  \draw[arr] (XT) -- (YT) node[midway, above] {$c$};
  \draw[arr] (XB) -- (YB) node[midway, below] {$c$};
  \draw[arr] (XT) -- (XB) node[midway, left]  {$e$};
  \draw[arr] (YT) -- (YB) node[midway, right] {$d$};
  \node[note] at (15,-10) {$e^2 = e$, \ $d^2 = d$};
  \node[note] at (15,-17) {$dc = ce$};
  % Middle: the functor sending the square to a morphism of images.
  \draw[mapsto] (44,11) -- (60,11) node[midway, above] {$\Phi_E$};
  % Right: the induced comparison between the idempotent images.
  \node[panel] at (95,31) {Comparison between images};
  \node (L) at (75,11)  {$(X,e)$};
  \node (R) at (115,11) {$(Y,d)$};
  \draw[arr] (L) -- (R) node[midway, above] {$dce$};
  \node[note] at (95,-10) {$d(dce)e = dce$};
\end{tikzpicture}

\caption{\textbf{From an idempotent change of a comparison to a comparison between the images.}
On the left we require $e^2=e$, $d^2=d$, and the commutativity $dc=ce$.
On the right the two ends are changed to the idempotent images, and the comparison between
them is $dce$.
In general $dce=c$ need not hold, so the round trip in the other direction returns to the
original by the isomorphism in the proof.}\label{fig:6.1}
\end{figure}

\begin{proof}

From $dc=ce$ and idempotence, $dce$ satisfies \eqref{eq:6.9}.
Let $a,b$ be the two ends of a morphism of the left-hand side. In addition to the original
commutative square, $e'ae=a$ and $d'bd=b$ hold.
Hence, regarding the same $a,b$ as Karoubi morphisms,
$b(dce)=(d'c'e')a$ holds.
Identities and composition are the same computation at the two ends, so this gives a
functor.

In the opposite direction, from a Karoubi morphism $u:(X,e)\to(Y,d)$ we take the underlying
morphism $u:X\to Y$ and the idempotent square $(e,d)$. The square commutes by $du=u=ue$.
For morphisms as well we take the underlying morphisms at the two ends.
Applying $\Phi_E$ after this returns to the original, since $due=u$.

In the round trip in the other direction, the comparison of $(c,e,d)$ changes to $dce$.
Placing the squares whose two ends are $(e,d)$ in both directions, we obtain an
isomorphism inside the idempotent completion.
The reason is that the two ends of both composites are $(e,d)$, which are the identity
morphisms there.
Naturality of this isomorphism also follows from the equations $e'ae=a$ and $d'bd=b$ at the
two ends of a morphism.

\end{proof}

\begin{proposition}[Compatibility with functors and projections]\label{prop:6.28}

A functor $F:E\to E'$ induces $\mathrm{Kar}(F)$ by $(X,e)\mapsto(FX,F(e))$ and its action
on morphisms.
This action and \eqref{eq:6.41} commute under a natural isomorphism.
In particular, the three-level projections induce functors
\begin{equation*}
\mathcal{M}(E_{\mathrm{geom}})
\longrightarrow\mathcal{M}(E_{\mathrm{core}})
\longrightarrow\mathcal{M}(B)
\tag{6.43}\label{eq:6.43}
\end{equation*}
and the projection obtained by following two levels agrees with the composite projection.

\end{proposition}

\begin{proof}

A functor preserves idempotence and \eqref{eq:6.9}.
The comparison obtained by applying the functor to \eqref{eq:6.42} is $F(dce)$, and the
comparison obtained by sending each component first is $F(d)F(c)F(e)$. The two are equal
because a functor preserves composites.
The endpoints and the components at the two ends of a morphism are the same as well, and we
obtain a natural isomorphism whose components are the identity morphisms at the endpoints.
Carrying out the same computation of components for the composite of two functors yields
the last claim.

\end{proof}

\begin{construction}[Invertible changes of comparisons]\label{cons:6.29}

Keeping all the objects of $\mathcal{M}(E)$ and taking as morphisms only the isomorphisms
between them yields the maximal groupoid.
We do not ask for invertibility of the comparison $c$ itself, which serves as an object.
What is required to be invertible are the changes at the two ends that match the
comparisons.
Indeed, an isomorphism of the arrow category is equivalent to a commutative square whose
two ends are isomorphisms.
One direction follows by evaluating at the two ends, and the other follows because the
inverse morphisms at the two ends form the inverse commutative square.

There are two ways of reading the generated comparison $\beta$ as an object.
They are $J(\beta):J(X)\to J(Y)$, the embedding of the original morphism, and
$\beta:(X,e)\to(Y,d)$, the morphism between the images of Theorem~\ref{thm:6.12}.
Even when the former is not invertible, the latter is an isomorphism.
By making these endpoints explicit, Chapter~\ref{chap:7} can specify ``which change
preserves which comparison.''

\end{construction}

\section{What remains in the image for lenses and protocols}\label{sec:6.8}

We examine whether reads, updates, and executions can be continued with the remaining
states after the states have been collected by an idempotent.
For fixed-view lenses and protocols with observations, the image of a semantics-preserving
idempotent can be realized inside each semantics.

\begin{proposition}[Splitting of the idempotents of fixed-view lenses]\label{prop:6.30}

The category $\mathrm{Lens}(V,v_0)$ of Definition~\ref{def:1.32} is idempotent complete.
That is, every idempotent semantics-preserving morphism $h:L\to L$ splits within the same
category.

\end{proposition}

\begin{proof}

By Propositions~\ref{prop:1.33} and~\ref{prop:1.34}, under the product decomposition
$C\cong V\times K_L$ of a lens $L$, an endomorphism can be written as
\[
h(v,k)=(v,t(k)),\qquad t:K_L\to K_L
\]
If $h^2=h$, then $t^2=t$, so we take
\[
K_t=\{k\in K_L\mid t(k)=k\},\qquad C_t=V\times K_t
\]
and take the read to be the first projection and the update to be
$((v,k),w)\mapsto(w,k)$.
Under the product decomposition, setting $r:C\to C_t$ to be $r(v,k)=(v,t(k))$ and
$i:C_t\to C$ to be the inclusion $i(v,k)=(v,k)$, both morphisms preserve reads and updates.
By idempotence $t(k)\in K_t$, and at fixed points $t(k)=k$, so we obtain
$ri=1_{C_t}$ and $ir=h$.
The reference fiber is in one-to-one correspondence with $K_t$, a subset of the finite set
$K_L$, so this is a splitting within the same category of lenses.

\end{proof}

We check the splitting of Proposition~\ref{prop:6.30} on a lens with four states.
We take the product lens of the two-element sets $V=R=\{0,1\}$, with states $C=V\times R$.
We take the reference view to be 0, the read to be $g(v,r)=v$, and the update to be
$p((v,r),w)=(w,r)$.
The map that aligns the second component to 0,
\begin{equation*}
h(v,r)=(v,0)
\tag{6.44}\label{eq:6.44}
\end{equation*}
defines an endomorphism of $\mathrm{Lens}(V,0)$. Indeed,
\begin{equation*}
gh=g,\qquad
h(p((v,r),w))=(w,0)=p(h(v,r),w),\qquad h^2=h.
\tag{6.45}\label{eq:6.45}
\end{equation*}
It forgets the distinction of the second component while preserving reads and updates.
The state set of the fixed points $C_0=V\times\{0\}$ is closed under updates and becomes a
lens by the same equations.
Taking $C\to C_0$ to be $h$ and the map in the opposite direction to be the inclusion, and
using the identity on views in both cases, we see that this idempotent splits within the
category of
lenses.

Hence the image obtained by applying Theorem~\ref{thm:6.12} to the category of lenses can,
in this example, be expressed as the actual lens $C_0\to V$.
Processing that uses only reads and updates can be continued on it, while processing that
asks for the original second component needs the original states.

In this way, for lenses over a fixed view, the image of a normalization can be realized
again as a lens.
For the normalization of cores treated in Theorem~\ref{thm:6.8}, what was required of an intermediate object was a core carrying the readings for all
the objects, together with its
morphisms.
Even in the same discussion of idempotents, the conclusion changes according to the
category in which the image is realized.

\begin{proposition}[Images of protocols preserving executions]\label{prop:6.31}

Fix the category of paths $\mathcal{C}_Q$ of Definition~\ref{def:1.36} and the observation
targets $O$, and take a realization $F:\mathcal{C}_Q\to{\mathbf{Set}}$ and an observation
$o_F:F\Rightarrow O$.
The state set at each vertex is taken to be finite.
Suppose a semantics-preserving adapter $n:F\Rightarrow F$ satisfies $n^2=n$.
In particular, we are assuming naturality with respect to all executions and preservation
of observations, $o_Fn=o_F$.
At each vertex we set
\begin{equation*}
F_n(x)=\{s\in F(x)\mid n_x(s)=s\}
\tag{6.46}\label{eq:6.46}
\end{equation*}
Restricting the executions and the observations then yields a realization $(F_n,o_{F_n})$
of the same protocol.
The maps $r_x(s)=n_x(s)$ and the inclusions $i_x$ are adapters preserving executions and
observations, and
\begin{equation*}
ri=1_{F_n},\qquad ir=n.
\tag{6.47}\label{eq:6.47}
\end{equation*}

\end{proposition}

\begin{proof}

For an execution $a:x\to y$ and a fixed point $s\in F_n(x)$, naturality gives
\begin{equation*}
n_y(F(a)(s))=F(a)(n_x(s))=F(a)(s)
\tag{6.48}\label{eq:6.48}
\end{equation*}
so the result after the execution again belongs to the fixed points.
For identities and composition it suffices to restrict the functor laws of $F$.
Naturality of $r$ is naturality of $n$, and naturality of $i$ is by the definition of the
restriction.
Setting the observation to be $o_{F_n}=o_Fi$, we have $o_{F_n}r=o_Fn=o_F$, so $r$ preserves
observations as well.
Each set of fixed points is a subset of the original finite set, so finiteness is preserved
as well.
Equation~\eqref{eq:6.47} follows from the definition of fixed points and idempotence.

\end{proof}

Having an idempotent map of states at each vertex alone does not give \eqref{eq:6.48}.
For example, realize a protocol with one vertex, one operation, and no declared relations
by two states and a constant observation.
Take the operation to be the exchange of 0 and 1, and the normalization to be the constant
map to 0.
Then the execution from the fixed point 0 goes out to 1.
In this case the normalization of states and the operation do not commute.

Equations \eqref{eq:6.45} for lenses and \eqref{eq:6.48} for protocols give the condition
under which operations can be continued with only the information retained after the
normalization.
Given a semantics-preserving isomorphism $\alpha$ and such an idempotent adapter $d$, the
morphism $d\alpha$ of Theorem~\ref{thm:6.12} matches the images at the two ends by an
isomorphism.
The reason why Chapter~\ref{chap:1} kept the names and the actions of operations separately
also appears here.
By checking the correspondence of the selected operations in addition to the agreement of
the actions on states and configurations, one can make clear which meaning the
normalization preserves.

\section*{Summary of the Chapter}

In this chapter we analyzed the generated comparison through its normalization factor and
clarified its invertibility and the structure that remains in the image.

\begin{itemize}
\item \textbf{Normalization and the information retained.} The map to the selected object
for each configuration is idempotent, and the fixed points are in one-to-one correspondence
with the configurations.
A readout unchanged by normalization factors uniquely through the configurations.
Under the conditions preserving operations, residuals, invariants, and signatures, this map
becomes an idempotent of full geometries (\S\S\ref{sec:6.1}--\ref{sec:6.2},
\S\ref{sec:6.4}).
\item \textbf{Invertibility and image of the generated comparison.} The generated
comparison, obtained by composing an invertible canonical comparison with a normalization,
is invertible if and only if the normalization factor is the identity.
The same comparison morphism gives an isomorphism between the images in the idempotent
completion even when the normalization is not the identity, and is compatible with the
projection to cores (\S\S\ref{sec:6.3}--\ref{sec:6.4}).
\item \textbf{Normalized changes.} The objects and morphisms after normalization form a category, and the normalization functor obtained from the original changes is full.
Whether the order of an original change and the normalization can be interchanged is
determined by the correspondence of the operations themselves, and it can fail to hold even
when only the actions on objects and configurations agree
(\S\S\ref{sec:6.5}--\ref{sec:6.6}).
\item \textbf{The category of comparisons themselves.} Combining the arrow category with
idempotent completion, comparisons and their changes can be handled.
The interchange of the two constructions is compatible with the three-level projections of
extraction, core, and geometry (\S\ref{sec:6.7}).
\item \textbf{Images on which operations can be continued.} The idempotents of fixed-view
lenses split within the same category.
For protocols as well, executions and observations are restricted to the image of an
idempotent adapter that preserves executions and observations (\S\ref{sec:6.8}).
\end{itemize}

For the design review of the opening, we select a standard representation that preserves
the operations, the dependencies, and the Laws of order cancellation.
Even when the two proposals after normalization can be compared by an isomorphism, the
original class structure and the auxiliary representations need not be recoverable.
The classification of this chapter answers separately what remains in the image and whether
one can return to the original representation.
To move a further design change across the normalization, one checks that the
correspondence of the operations commutes with the normalization.

In the next chapter we use the category of comparisons constructed here to classify the
changes that preserve a single comparison.
The next question is how far the changes that look the same after normalization can be
distinguished on the original objects.
